
\medskip
\phantomsection\label{prop:sga3-VIB-4.2}%
\noindent\textbf{Proposition 4.2.---}
\emph{Let \(\pi:G\to S\) be an \(S\)-scheme locally of finite
presentation, equipped with an \(S\)-section \(\epsilon\), and
satisfying the following two conditions \textup{(which hold when
\(G\) is an \(S\)-group, by \hyperref[lem:sga3-VIB-1.5]{1.5} and VI\(_A\), \SGARefLink{cor:sga3-VIA-2.4.1}{2.4.1})}:}
\begin{enumerate}
\item[\textup{\SGAInlineRefTarget{prop:sga3-VIB-4.2:item:a}(a)}] \emph{For every \(s\in S\) and every
  \(x\in G_s\), one has}
  \[
    \dim G_s=\dim_x G_s
  \]
  \emph{\textup{(or equivalently \textup{(\hyperref[lem:sga3-VIB-1.5]{1.5})}, for every
  \(s\in S\), all the irreducible components of \(G_s\) have the
  same dimension)}.}
\item[\textup{(b)}] \emph{For every \(s\in S\), if \(\underline{G_s^0}\)
  denotes the connected component of \(G_s\) containing
  \(\epsilon(s)\), then \(\underline{G_s^0}\) is geometrically irreducible.}
\end{enumerate}
\emph{Let \(s\in S\). The following conditions are equivalent:}
\begin{enumerate}
\item[\textup{\SGAInlineRefTarget{prop:sga3-VIB-4.2:item:i}(i)}] \emph{\(G\) is universally open over \(S\) at
  the points of \(\underline{G_s^0}\).}
\item[\textup{(i bis)}] \emph{\(G\) is universally open over \(S\)
  at every point of a neighborhood of \(\epsilon(s)\) in
  \(\underline{G_s^0}\).}%
\item[\textup{\SGAInlineRefTarget{prop:sga3-VIB-4.2:item:ii}(ii)}] \emph{The function}
  \[
    t\longmapsto\dim G_t
  \]
  \emph{is constant in a neighborhood of \(s\) in \(S\).}
\item[\textup{(iii)}]\phantomsection\label{concl:sga3-VIB-4.2-iii} \emph{\(\underline{G^0}\) is ``universally open over
  \(S\) at the points of \(\underline{G_s^0}\),'' that is, given
  \(S'\to S\), \(s'\in {S'}_s\), \(g\in\underline{G_{s'}^0}\), and an open
  neighborhood \(V\) of \(g\) in \(G'=G_{S'}\), then}
  \[
    \pi\bigl(V\cap\underline{{G'}^0}\bigr)
  \]
  \emph{is an open neighborhood of \(s'\) in \(S'\).}%
  \footnote{\label{note:sga3-VIB-29}Editor's note: we have added
  condition \textup{\SGARefLink{concl:sga3-VIB-4.2-iii}{(iii)}}, pointed out by O.\ Gabber; it will be
  useful later \textup{(\SGARefLink{thm:sga3-VIB-5.7}{5.7})}.}
\end{enumerate}

Clearly (i) implies (i bis). By EGA~IV\(_3\), 14.3.3.1(ii), the set
of points of \(\underline{G_s^0}\) at which \(\pi_G\) is universally open is
closed in \(\underline{G_s^0}\). Since \(\underline{G_s^0}\) is irreducible,
it follows that
(i bis) implies \SGARefLink{prop:sga3-VIB-4.2:item:i}{(i)}.

We show that \SGARefLink{prop:sga3-VIB-4.2:item:i}{(i)} implies \SGARefLink{prop:sga3-VIB-4.2:item:ii}{(ii)}. Let \(T\) be the set of
\(t\in S\) such that
\[
  \dim G_t=\dim G_s.
\]
By 4.1, \(T\) is locally constructible; hence, by
EGA~IV\(_1\), 1.10.1, to show that \(T\) is a neighborhood of
\(s\), it is enough to show that every generalization \(s'\) of
\(s\) belongs to \(T\).

Let \(x\) be the generic point of \(\underline{G_s^0}\), and let \(U\) be an
open subset of \(G\) containing \(x\). Since \(\pi_G\) is
universally open at \(x\), EGA~IV\(_3\), 14.3.13 gives, for every
generalization \(s'\) of \(s\),
\[
  \dim\bigl(U\cap G_{s'}\bigr)
    \geq \dim_x\bigl(U\cap G_s\bigr).
\]
In view of hypothesis \SGARefLink{prop:sga3-VIB-4.2:item:a}{(a)}, this implies
\[
  \dim G_{s'}\geq\dim G_s.
\]
On the other hand, by \hyperref[prop:sga3-VIB-4.1]{4.1} the function
\(s\mapsto\dim G_s\) is upper semicontinuous, so that
\[
  \dim G_{s'}\leq\dim G_s.
\]
Thus \(s'\in T\), proving that \SGARefLink{prop:sga3-VIB-4.1:item:i}{(i)} implies \SGARefLink{prop:sga3-VIB-4.1:item:ii}{(ii)}.

It is clear that \SGARefLink{concl:sga3-VIB-4.2-iii}{(iii)} implies \SGARefLink{prop:sga3-VIB-4.1:item:i}{(i)}; we prove that \SGARefLink{prop:sga3-VIB-4.2:item:ii}{(ii)} implies
\SGARefLink{concl:sga3-VIB-4.2-iii}{(iii)}. The dimension of the fibers is unchanged by extension of
the base field, and the formation of \(\underline{G^0}\) commutes
with base change (see the proof of \hyperref[prop:sga3-VIB-3.3]{3.3}). We may therefore assume
that \(S'=S\) and \(s'=s\). Moreover, every open subset \(V\) of
\(G\) meeting \(\underline{G_s^0}\) contains the generic point \(\eta\) of
\(\underline{G_s^0}\), so we may assume that \(g=\eta\).

We may further suppose that \(S\) is affine. Let \(U\) be an
affine open neighborhood of \(\epsilon(s)\); it is then of finite
presentation over \(S\). Replacing \(S\) by \(\epsilon^{-1}(U)\),
and then \(U\) by \(U\cap\pi^{-1}(S)\), we reduce to the case in
which
\[
  \pi:U\longrightarrow S
\]
is of finite presentation and admits
\(\epsilon:S\to U\) as a section. By EGA~IV\(_3\), 9.7.12,
\(\underline{G^0}\cap U\) is constructible in \(U\). Replacing
\(V\) by an affine open subset contained in \(V\cap U\), we find
that \(\underline{G^0}\cap V\) is constructible in \(V\). Since
\(\pi:V\to S\) is of finite presentation, Chevalley's
constructibility theorem (EGA~IV\(_1\), 1.8.4) shows that
\[
  \pi\bigl(\underline{G^0}\cap V\bigr)
\]
is locally constructible in \(S\).

Consequently, by \emph{loc.\ cit.}, 1.10.1, to show that
\(\pi(\underline{G^0}\cap V)\) is an open neighborhood of \(s\),
it is enough to show that, for every generalization \(t\) of
\(s\), there is a generalization \(\xi\) of \(\eta\) belonging to
\(\underline{G^0}\), and hence to
\(\underline{G^0}\cap V\). The generic point \(\xi\) of
\(\underline{G_t^0}\)
is such a generalization of \(\eta\). Indeed, \(\epsilon(s)\)
belongs to the closure \(X\) of \(\{\xi\}\) in \(G\); hence,
by Chevalley's semicontinuity theorem (EGA~IV\(_3\), 13.1.3),
\[
  \dim_{\epsilon(s)}X_s\geq\dim_\xi X_t.
\]
On the other hand, hypothesis \SGARefLink{prop:sga3-VIB-4.2:item:ii}{(ii)} implies
\[
  \dim_\xi\underline{G_t^0}=\dim_{\epsilon(s)}G_s.
\]
It follows that one of the irreducible components of \(X_s\)
containing \(\epsilon(s)\) is equal to \(\underline{G_s^0}\), whence
\[
  \eta\in\overline{\{\xi\}}.
\]
This proves that \SGARefLink{prop:sga3-VIB-4.2:item:ii}{(ii)} implies \SGARefLink{concl:sga3-VIB-4.2-iii}{(iii)}, and hence proves the
proposition.

\footnote{\label{note:sga3-VIB-30}Editor's note: the preceding
argument was communicated to us by O.\ Gabber. We have also
retained the proof of the implication \textup{\SGARefLink{prop:sga3-VIB-4.2:item:ii}{(ii)}}
\(\Rightarrow\)\textup{\SGARefLink{prop:sga3-VIB-4.2:item:i}{(i)}} given in the original.}%
One can also prove that \SGARefLink{prop:sga3-VIB-4.2:item:ii}{(ii)} implies \SGARefLink{prop:sga3-VIB-4.2:item:i}{(i)} as follows. Since \(\pi\)
is locally of finite presentation, there exist an open subset
\(U\) of \(G\) containing \(\epsilon(s)\) and an open subset
\(V\) of \(S\) containing \(s\), such that \(\pi(U)\subset V\)
and the morphism
\[
  \pi':U\longrightarrow V
\]
induced by \(\pi\) is of finite presentation. Set
\[
  T=\epsilon^{-1}(U),
  \qquad
  W={\pi'}^{-1}(T)=U\cap\pi^{-1}(T).
\]
The morphism
\[
  \pi'':W\longrightarrow T
\]
induced by \(\pi'\) is of finite presentation and admits as a
section the morphism
\[
  \epsilon'':T\longrightarrow W
\]
induced by \(\epsilon\). Moreover, for every \(t\in T\), since
\(\underline{G_t^0}\) is irreducible,
\(W\cap\underline{G_t^0}\) is dense in \(\underline{G_t^0}\),
hence irreducible and therefore connected. It is consequently the
connected component of \({\pi''}^{-1}(t)\) containing
\(\epsilon''(t)\).

Since \(W\cap\underline{G_t^0}\) is a dense open subset of
\(\underline{G_t^0}\), \hyperref[lem:sga3-VIB-1.5]{1.5}
and EGA~IV\(_2\), 5.2.1 give
\[
  \dim(W\cap\underline{G_t^0})=\dim\underline{G_t^0}=\dim G_t.
\]
Thus the function
\[
  t\longmapsto\dim(W\cap\underline{G_t^0})
\]
is constant in a neighborhood of \(s\) in \(T\). Finally, we show
that \(W\cap\underline{G_t^0}\) is geometrically irreducible for every
\(t\in T\). Let \(K\) be an extension of \(\kappa(t)\). Then
\[
  (W\cap\underline{G_t^0})\otimes_{\kappa(t)}K
    =
  \bigl(W\otimes_{\kappa(t)}K\bigr)
    \cap
  \bigl(\underline{G_t^0}\otimes_{\kappa(t)}K\bigr)
\]
is a nonempty open subset of
\(\underline{G_t^0}\otimes_{\kappa(t)}K\), and is therefore irreducible
because \(\underline{G_t^0}\otimes_{\kappa(t)}K\) is irreducible.

We are now in the situation of EGA~IV\(_3\), 15.6.6(ii), which
asserts that \(\pi''\), and hence \(\pi\), is universally open at
the points of \(W\cap\underline{G_s^0}\). But, by
EGA~IV\(_3\), 14.3.3.1(ii),
the subset of \(G_s\) consisting of the points at which \(\pi\)
is universally open is closed in \(G_s\). Since it contains
\(W\cap\underline{G_s^0}\), it therefore contains its closure
\(\underline{G_s^0}\).
